% ANNEXE ST
% ==============================================================
\sectionintro{Annexe A · TwinCAT 3}{Types et interface du bloc}{Le code de la version finale est regroupé dans le fichier \texttt{FB\_Carrousel.st} joint au rapport.}

\noindent
\begin{minipage}[t]{0.47\textwidth}\vspace{0pt}
  \subsection{Énumérations}
\begin{lstlisting}[style=stcode]
{attribute 'qualified_only'}
{attribute 'strict'}
{attribute 'to_string'}
TYPE E_ModeCarrousel :
(
    INIT   := 0,
    AUTO   := 1,
    DEFAUT := 2
);
END_TYPE

{attribute 'qualified_only'}
{attribute 'strict'}
{attribute 'to_string'}
TYPE E_EtatCycle :
(
    ATTENTE       := 0,
    COMPTAGE      := 1,
    BAC_PLEIN     := 2,
    VERIN_SORTIE  := 3,
    VERIN_RENTREE := 4,
    BAC_SUIVANT   := 5,
    ATTENTE_BAC   := 6
);
END_TYPE
\end{lstlisting}

  \begin{notebox}{Lecture en ligne}
    Les énumérations rendent le mode et l'état courant directement compréhensibles dans TwinCAT, sans devoir interpréter des numéros d'étape.
  \end{notebox}
\end{minipage}\hfill
\begin{minipage}[t]{0.49\textwidth}\vspace{0pt}
  \subsection{Interface de \texttt{FB\_Carrousel}}
\begin{lstlisting}[style=stcode]
FUNCTION_BLOCK FB_Carrousel
VAR_INPUT
    bLaser    : BOOL; // detection piece
    bZero     : BOOL; // origine carrousel
    nPieceBac : INT;  // consigne / bac
    tVerin    : TIME; // temps mouvement
    aResetBac : ARRAY[1..8] OF BOOL;
END_VAR

VAR_OUTPUT
    bVerinAv : BOOL; // cmd avant
    bVerinAr : BOOL; // cmd arriere
    nBac     : INT;  // bac en position
    bDefaut  : BOOL; // defaut carrousel
END_VAR

VAR
    eMode      : E_ModeCarrousel
                 := E_ModeCarrousel.INIT;
    eEtatCycle : E_EtatCycle
                 := E_EtatCycle.ATTENTE;
    aPieceBac  : ARRAY[1..8] OF INT;
    nImp       : INT := 0;
    fbLaser      : R_TRIG;
    fbTempoVerin : TON;
    aTempoReset  : ARRAY[1..8] OF TON;
    bResetGlobal : BOOL;
    i, j         : INT;
END_VAR
\end{lstlisting}

  \begin{primarycard}
    \textbf{Organisation du source joint}\par\small
    Le fichier regroupe les deux DUTs, le bloc fonction et un exemple de \texttt{MAIN}. Dans le projet XAE, ces éléments restent répartis dans leurs objets TwinCAT respectifs.
  \end{primarycard}
\end{minipage}

